\documentclass[10pt,
% a4paper,
%twocolumn,
fleqn,
%landscape, 
% papersize,
dvipdfmx,
uplatex
]{jsarticle}



\def\maru#1{\textcircled{\scriptsize#1}}%丸囲み番号

% \RequirePackage[2020/09/30]{platexrelease}

%太字設定
\usepackage[deluxe]{otf}

\usepackage{emathEy}

\usepackage[g]{esvect}

%大きな文字
\usepackage{fix-cm}

%定理環境
\usepackage{emathThm}
%\theoremstyle{boxed}
\theorembodyfont{\normalfont}
\newtheorem{Question}{問題}[subsection]
\newtheorem{Q}{}[subsection]
\newtheorem{question}[Question]{}
\newtheorem{quuestion}{}[subsection]

%セクション，大問番号のデザイン
\renewcommand{\labelenumi}{(\arabic{enumi})}
\renewcommand{\theenumii}{\alph{enumii})}
\renewcommand{\thesection}{第\arabic{section}章}

%用紙サイズの詳細設定
\usepackage{bxpapersize}
\papersizesetup{size={80mm,45mm}}
\usepackage[top=0.7zw,bottom=0truemm,left=3truemm,right=133truemm]{geometry}
\usepackage[dvipdfmx]{graphicx}

%余白など
\usepackage{setspace} % 行間
\setlength{\mathindent}{1zw}
\setlength\parindent{0pt}


%色カラーに関する設定
\usepackage{color}
\definecolor{shiro}{rgb}{0.95703125,0.87109375,0.7421875}
\definecolor{kin}{rgb}{0.95703125,0.87109375,0.7421875}
\definecolor{orange}{rgb}{1,0.7,0.2}
\definecolor{bradorange}{rgb}{1,0.5,0}
\definecolor{pink}{rgb}{0.9176,0.5686,0.5960}
\definecolor{mizu}{rgb}{0.6156,0.8,0.9955}
\color{kin}
% \pagecolor{hukamido}

\usepackage{at}%図の配置
% \usepackage{wallpaper}

\begin{document}



\at(0cm,0cm){\includegraphics[width=8cm,bb=0 0 1920 1080]{./youtube/thumbnails/templates/smart_background/場合の数.jpeg}}
{\color{orange}\bf\boldmath\LARGE\underline{並ぶ順番についての条件}}\vspace{0.3zw}

\normalsize 
\bf\boldmath 問．N,\;I,\;S,\;H,\;I,\;O,\;J,\;Iの$8$文字を一列に並べるとき，$4$つの子音

\huge
\vspace{-0.2zw}
\hspace{0.1zw} N,\;S,\;H,\;Jが\vspace{-0.2zw}\\
\hfill この順に並ぶもの\hspace{0.1zw} 

\normalsize
\hfill は通りあるか．
\at(7.0cm,0.2cm){\small\color{bradorange}$\overset{\text{場合の数}}{\text{典型}}$}


\newpage



\at(0cm,0cm){\includegraphics[width=8cm,bb=0 0 1920 1080]{./youtube/thumbnails/templates/smart_background/場合の数.jpeg}}
{\color{orange}\bf\boldmath\LARGE\underline{円卓で向かい合う座り方}}\vspace{0.3zw}

\normalsize 
\bf\boldmath 問．大人$2$人と子ども$4$人が円卓に座るとき，

\Huge
\vspace{-0.2zw}
\hspace{0.05zw} 大人$2$人が向かい\vspace{-0.2zw}\\
\hfill 合うような座り方\hspace{0.1zw} 
\vspace{0.1zw}

\normalsize
\hfill は何通りあるか．
\at(7.0cm,0.2cm){\small\color{bradorange}$\overset{\text{場合の数}}{\text{典型}}$}


\newpage



\at(0cm,0cm){\includegraphics[width=8cm,bb=0 0 1920 1080]{./youtube/thumbnails/templates/smart_background/場合の数.jpeg}}
{\color{orange}\bf\boldmath\Large\underline{隣り合う$\cdot$隣り合わない並び方}}\vspace{0.3zw}

\normalsize 
\bf\boldmath 問．男子$5$人と女子$3$人を一列に並べるとき，次のような並べ方は何通りあるか．

\huge
% \vspace{-0.2zw}
(1)女子$3$人が隣り合う

(2)女子が隣り合わない

\at(7.0cm,0.2cm){\small\color{bradorange}$\overset{\text{場合の数}}{\text{典型}}$}


\newpage



\at(0cm,0cm){\includegraphics[width=8cm,bb=0 0 1920 1080]{./youtube/thumbnails/templates/smart_background/場合の数.jpeg}}
{\color{orange}\bf\boldmath\large\underline{頂点を結んでできる三角形の個数}}\vspace{0.1zw}

\normalsize 
\bf\boldmath 問．九角形の

\huge
\vspace{-0.2zw}
\hspace{0.1zw} $9$個の頂点の$3$つを\\
\hfill 結んでできる三角形 \hspace{0.1zw}

\normalsize
\vspace{-0.2zw}
のうち，\vspace{-0.2zw}\\
\hfill 九角形と辺を共有しないものはいくつあるか．
\at(7.0cm,0.2cm){\small\color{bradorange}$\overset{\text{場合の数}}{\text{典型}}$}


\newpage



\at(0cm,0cm){\includegraphics[width=8cm,bb=0 0 1920 1080]{./youtube/thumbnails/templates/smart_background/場合の数.jpeg}}
{\color{orange}\bf\boldmath\huge\underline{対角線の本数}}\vspace{0.3zw}

\large
\bf\boldmath 問．

\HUGE
\vspace{-0.6zw}
\hspace{0.00zw}
九角形の対角線\\
\hfill は何本あるか．
\at(7.0cm,0.2cm){\small\color{bradorange}$\overset{\text{場合の数}}{\text{典型}}$}


\newpage



\at(0cm,0cm){\includegraphics[width=8cm,bb=0 0 1920 1080]{./youtube/thumbnails/templates/smart_background/場合の数.jpeg}}
{\color{orange}\bf\boldmath\Large\underline{同じものを含む円順列 Lv.2 }}\vspace{0.3zw}

\huge 
\bf\boldmath 問．赤玉$4$個，白玉$2$個，黒玉$2$個を円形に並べる方法は何通りか．
\at(7.0cm,0.2cm){\small\color{bradorange}$\overset{\text{場合の数}}{\text{典型}}$}


\newpage



\at(0cm,0cm){\includegraphics[width=8cm,bb=0 0 1920 1080]{./youtube/thumbnails/templates/smart_background/場合の数.jpeg}}
{\color{orange}\bf\boldmath\Large\underline{同じものを含む円順列 Lv.1 }}\vspace{0.3zw}

\huge 
\bf\boldmath 問．赤玉$4$個，白玉$3$個，黒玉$1$個を円形に並べる方法は何通りか．
\at(7.0cm,0.2cm){\small\color{bradorange}$\overset{\text{場合の数}}{\text{典型}}$}


\newpage



\at(0cm,0cm){\includegraphics[width=8cm,bb=0 0 1920 1080]{./youtube/thumbnails/templates/smart_background/場合の数.jpeg}}
{\color{orange}\bf\boldmath\LARGE\underline{重複組合せ Lv.2 }}\vspace{0.3zw}

\normalsize
\bf\boldmath 問．${10}$個のミカンをA,\;B,\;Cの$3$人配る．

\huge
\vspace{-0.1zw}
\hspace{0.1zw} $1$つももらわない人が\\
\hfill いても良いとき\hspace{0.1zw} 

\normalsize
\vspace{0.0zw}
\hfill の配り方は何通りあるか．

\at(7.0cm,0.2cm){\small\color{bradorange}$\overset{\text{場合の数}}{\text{典型}}$}


\newpage



\at(0cm,0cm){\includegraphics[width=8cm,bb=0 0 1920 1080]{./youtube/thumbnails/templates/smart_background/場合の数.jpeg}}
{\color{orange}\bf\boldmath\LARGE\underline{重複組合せ Lv.1 }}\vspace{0.3zw}


\normalsize
\bf\boldmath 問．${10}$個のミカンをA,\;B,\;Cの$3$人配る．

\huge
\vspace{-0.1zw}
\hspace{0.1zw} どの人も少なくとも\\
\hfill$1$個はもらうとき\hspace{0.1zw} 

\normalsize
\vspace{0.1zw}
\hfill の配り方は何通りあるか．

\at(7.0cm,0.2cm){\small\color{bradorange}$\overset{\text{場合の数}}{\text{典型}}$}


\newpage



\at(0cm,0cm){\includegraphics[width=8cm,bb=0 0 1920 1080]{./youtube/thumbnails/templates/smart_background/場合の数.jpeg}}
{\color{orange}\bf\boldmath\huge\underline{組分け Lv.3 }}\vspace{0.3zw}

\large 
\bf\boldmath 問．${12}$人を

\Huge
\vspace{-0.3zw}
\hspace{0.1zw} $4$人，$4$人，$4$人の\vspace{-0.2zw}\\
\hfill $3$組に分ける方法\hspace{0.1zw} 

\large
\vspace{0.1zw}
\hfill は何通りか．
\at(7.0cm,0.2cm){\small\color{bradorange}$\overset{\text{場合の数}}{\text{典型}}$}


\newpage



\at(0cm,0cm){\includegraphics[width=8cm,bb=0 0 1920 1080]{./youtube/thumbnails/templates/smart_background/場合の数.jpeg}}
{\color{orange}\bf\boldmath\huge\underline{組分け Lv.2 }}\vspace{0.3zw}

\large 
\bf\boldmath 問．${12}$人を

\Huge
\vspace{-0.3zw}
\hspace{0.1zw} $6$人，$6$人の\vspace{-0.2zw}\\
\hfill $2$組に分ける方法\hspace{0.1zw} 

\large
\vspace{0.1zw}
\hfill は何通りか．
\at(7.0cm,0.2cm){\small\color{bradorange}$\overset{\text{場合の数}}{\text{典型}}$}


\newpage



\at(0cm,0cm){\includegraphics[width=8cm,bb=0 0 1920 1080]{./youtube/thumbnails/templates/smart_background/場合の数.jpeg}}
{\color{orange}\bf\boldmath\huge\underline{組分け Lv.1 }}\vspace{0.3zw}

\large 
\bf\boldmath 問．${12}$人を

\Huge
\vspace{-0.3zw}
\hspace{0.1zw} $5$人，$4$人，$3$人の\vspace{-0.2zw}\\
\hfill $3$組に分ける方法\hspace{0.1zw} 

\large
\vspace{0.1zw}
\hfill は何通りか．
\at(7.0cm,0.2cm){\small\color{bradorange}$\overset{\text{場合の数}}{\text{典型}}$}


\newpage



\at(0cm,0cm){\includegraphics[width=8cm,bb=0 0 1920 1080]{./youtube/thumbnails/templates/smart_background/場合の数.jpeg}}
{\color{orange}\bf\boldmath\huge\underline{組分け Lv.0 }}\vspace{0.3zw}

\large 
\bf\boldmath 問．${12}$人を

\Huge
\vspace{-0.3zw}
\hspace{0.1zw} $8$人，$4$人の\vspace{-0.2zw}\\
\hfill $2$組に分ける方法\hspace{0.1zw} 

\large
\vspace{0.1zw}
\hfill は何通りか．
\at(7.0cm,0.2cm){\small\color{bradorange}$\overset{\text{場合の数}}{\text{典型}}$}


\newpage



\at(0cm,0cm){\includegraphics[width=8cm,bb=0 0 1920 1080]{./youtube/thumbnails/templates/smart_background/場合の数.jpeg}}
{\color{orange}\bf\boldmath\huge\underline{部屋分け Lv.2 }}\vspace{0.3zw}


\large 
\bf\boldmath 問．空室は作らないものとする．

\Huge
\vspace{-0.3zw}
\hspace{0.1zw} $6$人をA,\;B,\;Cの\vspace{-0.2zw}\\
\hfill$3$部屋に分ける方法\hspace{0.1zw} 

\large
\vspace{0.1zw}
\hfill は何通りか．
\at(7.0cm,0.2cm){\small\color{bradorange}$\overset{\text{場合の数}}{\text{典型}}$}


\newpage



\at(0cm,0cm){\includegraphics[width=8cm,bb=0 0 1920 1080]{./youtube/thumbnails/templates/smart_background/場合の数.jpeg}}
{\color{orange}\bf\boldmath\huge\underline{部屋分け Lv.1 }}\vspace{0.3zw}

\large 
\bf\boldmath 問．空室は作らないものとする．

\Huge
\vspace{-0.3zw}
\hspace{0.1zw} $6$人をA,\;Bの\vspace{-0.2zw}\\
\hfill$2$部屋に分ける方法\hspace{0.1zw} 

\large
\vspace{0.1zw}
\hfill は何通りか．
\at(7.0cm,0.2cm){\small\color{bradorange}$\overset{\text{場合の数}}{\text{典型}}$}


\end{document}

